\begin{lemma}[Bijectivity of Fourier transform]
\label{bij}
\lean{fourier_surjection_on_schwartz}
% at /home/runner/work/trough/trough/library/PrimeNumberTheoremAnd/Wiener.lean:1828.0-1846.48
\leanok
The Fourier transform is a bijection on the Schwartz class. [Note: only surjectivity is
  actually used.]
\end{lemma}
\begin{proof}
\leanok
This is a standard result in Fourier analysis.
  It can be proved here by appealing to Mellin inversion, Theorem \ref{MellinInversion}.
  In particular, given $f$ in the Schwartz class, let
  $F : \R_+ \to \C : x \mapsto f(\log x)$ be a function in the ``Mellin space''; then the
  Mellin transform of $F$ on the imaginary axis $s=it$ is the Fourier transform of $f$.
  The Mellin inversion theorem gives Fourier inversion.
\end{proof}
